在平面内运动的质量 \(m>0\) 的质点受到吸引中心力 \(F_r=-k\,r^{p-1}\)，其中 \(k>0\)、实指数 \(p\in\mathbb R\)，角动量 \(L\ne0\)。当 \(p\ne0\) 时势能取 \(U(r)=kr^p/p\)；当 \(p=0\) 时定义 \(U(r)=k\ln(r/r_*)\)，其中 \(r_*>0\) 是任意参考长度。
(1) 写出径向有效势并由此求半径 \(r_0\) 的圆轨道所需角动量；(2) 严格判断圆轨道对小径向扰动稳定的 \(p\) 取值范围；(3) 在稳定范围内求径向振动与角向频率之比、相邻近心点间的方位角，判断对 \(p=-1\) 与 \(p=2\) 的闭合轨道结果；(4) 计算圆轨道的比动能及向心力平衡，说明为何 \(p\to0\) 的对数势极限是连续的。
